% GENERATED FILE. Edit canonical lessons, then run npm run generate:books.
% canonical-source-hash: fnv1a64:a471aa8f4bb72c95
% canonical-lessons: TE-C224-garite, TE-C224-cetta, TE-C224-nela, TE-C224-parupu, TE-C224-pindi

\chapter{Ladle, Rubbish, Floor, Mattress, Flour}
\label{ch:te-ladle-rubbish-floor-mattress-flour}
\hlchaptermodality{224}

\begin{chapteropening}
\textbf{By the end of this chapter:} \emph{I can say a ladle, rubbish, the floor, a mattress and flour.}

Complete the last lesson of chapter 224: I can say a ladle, rubbish, the floor, a mattress and flour.
\end{chapteropening}

\section[\texorpdfstring{\te{గరిటె} (gariṭe)}{గరిటె (gariṭe)}]{\te{గరిటె} (gariṭe) — a ladle}

\label{lesson:TE-C224-garite}



\begin{quote}
Before the new one: say the Telugu for messy, then the Telugu for rude.
\end{quote}

\subsection*{You'll want to know: \te{గరిటె}}
\textbf{\te{గరిటె}} — \emph{gariṭe} — \textquotedblleft{}a ladle\textquotedblright{}.

\begin{grammarlens}[title={What you've built}]
One more word for this chapter.
\end{grammarlens}

\subsection*{Your turn}
\begin{itemize}
\raggedright
  \item \emph{Say it:} \emph{gariṭe}
  \item \emph{Say it:} \emph{gariṭe}, once more
  \item \emph{Say it:} say \emph{gariṭe}
  \item \emph{Recall:} say the Telugu for messy, then the Telugu for rude, then say \emph{gariṭe} again
\end{itemize}

\begin{tcolorbox}[breakable,colback=teal!4,colframe=teal!35!black,title={Before you move on}]
What does \te{గరిటె} mean? (\textquotedblleft{}A ladle\textquotedblright{}.) Say it once more.
\end{tcolorbox}
\section[\texorpdfstring{\te{చెత్త} (cetta)}{చెత్త (cetta)}]{\te{చెత్త} (cetta) — rubbish, garbage}

\label{lesson:TE-C224-cetta}



\begin{quote}
Before the new one: say the Telugu for rude, then the Telugu for a ladle.
\end{quote}

\subsection*{You'll want to know: \te{చెత్త}}
\textbf{\te{చెత్త}} — \emph{cetta} — \textquotedblleft{}rubbish, garbage\textquotedblright{}.

\begin{grammarlens}[title={What you've built}]
One more word for this chapter.
\end{grammarlens}

\subsection*{Your turn}
\begin{itemize}
\raggedright
  \item \emph{Say it:} \emph{cetta}
  \item \emph{Say it:} \emph{cetta}, once more
  \item \emph{Say it:} say \emph{cetta}
  \item \emph{Recall:} say the Telugu for rude, then the Telugu for a ladle, then say \emph{cetta} again
\end{itemize}

\begin{tcolorbox}[breakable,colback=teal!4,colframe=teal!35!black,title={Before you move on}]
What does \te{చెత్త} mean? (\textquotedblleft{}Rubbish, garbage\textquotedblright{}.) Say it once more.
\end{tcolorbox}
\section[\texorpdfstring{\te{నేల} (nēla)}{నేల (nēla)}]{\te{నేల} (nēla) — the floor, the ground}

\label{lesson:TE-C224-nela}



\begin{quote}
Before the new one: say the Telugu for a ladle, then the Telugu for rubbish.
\end{quote}

\subsection*{You'll want to know: \te{నేల}}
\textbf{\te{నేల}} — \emph{nēla} — \textquotedblleft{}the floor, the ground\textquotedblright{}.

\begin{grammarlens}[title={What you've built}]
One more word for this chapter.
\end{grammarlens}

\subsection*{Your turn}
\begin{itemize}
\raggedright
  \item \emph{Say it:} \emph{nēla}
  \item \emph{Say it:} \emph{nēla}, once more
  \item \emph{Say it:} say \emph{nēla}
  \item \emph{Recall:} say the Telugu for a ladle, then the Telugu for rubbish, then say \emph{nēla} again
\end{itemize}

\begin{tcolorbox}[breakable,colback=teal!4,colframe=teal!35!black,title={Before you move on}]
What does \te{నేల} mean? (\textquotedblleft{}The floor, the ground\textquotedblright{}.) Say it once more.
\end{tcolorbox}
\section[\texorpdfstring{\te{పరుపు} (parupu)}{పరుపు (parupu)}]{\te{పరుపు} (parupu) — a mattress}

\label{lesson:TE-C224-parupu}



\begin{quote}
Before the new one: say the Telugu for rubbish, then the Telugu for the floor.
\end{quote}

\subsection*{You'll want to know: \te{పరుపు}}
\textbf{\te{పరుపు}} — \emph{parupu} — \textquotedblleft{}a mattress\textquotedblright{}.

\begin{grammarlens}[title={What you've built}]
One more word for this chapter.
\end{grammarlens}

\subsection*{Your turn}
\begin{itemize}
\raggedright
  \item \emph{Say it:} \emph{parupu}
  \item \emph{Say it:} \emph{parupu}, once more
  \item \emph{Say it:} say \emph{parupu}
  \item \emph{Recall:} say the Telugu for rubbish, then the Telugu for the floor, then say \emph{parupu} again
\end{itemize}

\begin{tcolorbox}[breakable,colback=teal!4,colframe=teal!35!black,title={Before you move on}]
What does \te{పరుపు} mean? (\textquotedblleft{}A mattress\textquotedblright{}.) Say it once more.
\end{tcolorbox}
\section[\texorpdfstring{\te{పిండి} (piṇḍi)}{పిండి (piṇḍi)}]{\te{పిండి} (piṇḍi) — flour}

\label{lesson:TE-C224-pindi}



\begin{quote}
Before the new one: say the Telugu for the floor, then the Telugu for a mattress.
\end{quote}

\subsection*{You'll want to know: \te{పిండి}}
\textbf{\te{పిండి}} — \emph{piṇḍi} — \textquotedblleft{}flour\textquotedblright{}.

\begin{grammarlens}[title={What you've built}]
One more word for this chapter.
\end{grammarlens}

\subsection*{Your turn}
\begin{itemize}
\raggedright
  \item \emph{Say it:} \emph{piṇḍi}
  \item \emph{Say it:} \emph{piṇḍi}, once more
  \item \emph{Say it:} say \emph{piṇḍi}
  \item \emph{Recall:} say the Telugu for the floor, then the Telugu for a mattress, then say \emph{piṇḍi} again
\end{itemize}

\begin{tcolorbox}[breakable,colback=teal!4,colframe=teal!35!black,title={Before you move on}]
What does \te{పిండి} mean? (\textquotedblleft{}Flour\textquotedblright{}.) Say it once more.
\end{tcolorbox}
